\section{Problem Statement}
\label{sec:problem}

\subsection{Application Model}
\label{sec:problem-application}

The application is a chain of $M{\geq}1$
tasks. Each task $T_i$, $i{=}1, \ldots, M$, takes input data,
performs a computation, and generates output data, as a three-step
sequence. Each successful execution of $T_i$ produces output data,
any of which may be used as input by an execution of $T_{i+1}$. 
Data is always written to and read from stable
storage.  An ``in-situ" execution in which the output data from a parent
task remains in RAM and can be accessed by the child task without
any I/O would simply amount to merging these two
tasks into a single task.


Error in application data is tolerated:  $T_i$ offers $K_i{\geq}1$
execution modes for different trade-offs between
execution time and output error. 
$T_i$ takes $x{\in}\N$ bytes of input with an error level $y{\in}[1,
+\infty)$, where 1 denotes no error.  We define execution mode $j$ for
$T_i$, $j{=}1,\ldots, K_i$, by three functions:
\begin{itemize}
\item $t^i_j(x,y)$: The compute time (not including I/O); 
\item $d^i_j(x,y)$: The size in bytes of the output data; and 
\item $e^i_j(x,y)$: The error level in the output data.  
\end{itemize}
\noindent
In practice, these functions may be independent of some of
their inputs (see the example in Table~\ref{tab:compression_example}).

\begin{table}[hb!]
    \caption{Example execution modes for a data compression task. $x$ is the input data size in bytes
	 and $y$ is the input data error.}
    \label{tab:compression_example}
    \begin{center}
    \begin{tabular}{|l|l|l|l|}
    \hline
    Execution mode & $t(x, y)$ & $d(x,y)$ & $e(x,y)$\\
    \hline
    no compression & $0$ & $x$ & $y$\\
    slow lossless compression & $2{+}10 x$ & $x/10$ & $y$\\
    fast lossless compression & $1{+}5 x$  & $x/2$  & $y$\\
    fast lossy compression    & $1{+}2 x$  & $x/20$ & $3y/2$\\
    \hline
    \end{tabular}
    \end{center}
\end{table}

Error propagates through inter-task data dependencies and
accumulates end-to-end.  Many previous works on the scheduling of
real-time, approximate computing applications employ the ``monotone" model
by which a task can reduce the error in its input data by performing
additional computation~\cite{259428}.  The above model can account for this
error reduction behavior by defining the $t$ and $e$ functions
appropriately.  For the use cases that motivate this work, as that in
Section~\ref{sec:use-case}, there is no possibility of error reduction
(i.e., $e^i_j(x,y){\geq}y$).

The execution of the application must complete by a deadline, $D{>}0$. The
\emph{output error} is defined as the minimum error across the output data
of all successful executions of $T_{M}$ (the last task).  If no 
execution of $T_{M}$ has succeeded by the deadline, then the 
output error is set to $\efail$, a user-defined value that quantifies
the penalty for returning no usable output by the deadline. The value of
$\efail$ is set to be larger than any achievable output error by a successful
execution of $T_M$. 

\subsection{Platform and Fault Model}
\label{sec:problem-execution}

The platform comprises $N$ homogeneous
processing elements (PEs), that each can execute only one task at a time.
A PE could be a single core that executes a sequential program,
a multi-core host that executes a multi-threaded program, or a set of compute
nodes that execute a parallel program.
Each PE reads/writes task data on a parallel file system,
with read bandwidth $B_r$ and write bandwidth $B_w$. In
practice, PEs can experience contention when accessing the parallel
file system concurrently. While contention is part of our simulation
model (for experiments), our analytical model assumes no
contention so that analytical developments remain tractable.

Task executions on PEs experience fail-stop failures (failed task executions do not produce usable
output data).  Each PE is subject to independent failures with 
exponentially distributed inter-arrival times with 
rate $\lambda$ ($\lambda{=}0$ corresponds to a failure-free system).  
Exponentially distributed failures are commonplace in the fault-tolerance
literature due to their memoryless property. Although failures are not
necessarily memoryless in real-world settings, assuming exponentially
distributed failures remains sensible in practice~\cite{8049007}. We assume
a single rate $\lambda$ regardless of a task's execution mode.  Some works
consider different failure rates for different execution modes~\cite{10.1145/2442087.2442094,10.1007/s11227-021-03764-x,10.1109/TC.2003.1183950}.  Considering different failure
rates for different modes in this work would merely entail simple
modifications to the recursive formulas presented in
Section~\ref{sec:onehost:solution}.

Following a failure, the affected PE cannot resume computing for $R{\geq}
0$ seconds. This \emph{restart overhead} corresponds to acquiring a spare
host (due to a hardware failure), rebooting the host (due to an OS
crash), or simply restarting a task (due to a software bug).  As
commonly done in the literature, we assume a single restart overhead and do
not distinguish different types of failures. Here again, doing so would
entail straightforward (if tedious) modifications to the recursive
formulas presented in Section~\ref{sec:onehost:solution}.

A key technique to handle fail-stop failures is rollback/recovery, by
which a task is restarted if it fails~\cite{10.5555/2811302}.  A common
question is that of whether and when a task should
checkpoint its state to stable storage to avoid rolling back the entire
computation when a failure occurs~\cite{BAUTISTAGOMEZ2024315}.  In this
work checkpointing is built into the application because each task reads
its input and saves its output to stable storage.  Our model can easily
be extended to consider a task that implements its own checkpointing:
replace the task by a chain of tasks, each of which corresponds to the
execution between two checkpoints. All tasks in this chain, except the
first, would read their input from stable storage only when restarted
after a failure, which would require minor changes to the recursive
formulas presented in Section~\ref{sec:onehost:solution}.  Besides
rollback/recovery, another key technique is replication, by which backup
tasks are executed on different PEs to increase the chances that at least
one replica will succeed~\cite{10.5555/2811302}.  In this work, we 
use both rollback/recovery and replication.

%There are works that say "checkpoint more the closer to the deadline": we
%don't do that. cite? (ref [68] in 3589950.pdf) Discuss in related work?
%Future work? Ignore?


\subsection{Scheduling Problem}
\label{sec:problem-scheduling}

We consider the problem of \emph{minimizing the expected output error} for
a single invocation of the application, released at time zero and subject
to an end-to-end deadline $D$.  As seen in our motivating use case in
Section~\ref{sec:use-case}, the application is intended to be invoked
periodically. But successive invocations do not overlap, no unfinished work
carries over beyond a deadline, and all PEs are available at the beginning
of each invocation.  Each invocation can thus be scheduled independently
using the same strategy. In other words, only the deployment of the
application is periodic, not its scheduling.

The scheduling problem is offline because application and platform specifications 
are known in advance.  But because
failures are stochastic, a scheduling strategy should be
adaptive and make decisions that account for execution events at runtime.
Therefore, we consider scheduling strategies that are invoked at time zero
and also in response to task completions, task failures, and PEs becoming
available.  Scheduling decisions can consist of aborting current task executions
(making PEs immediately available) or selecting which tasks to execute with
which execution modes.  We define an \emph{optimal strategy} as one in
which each decision minimizes the expected output error given the current
state of the application execution and the time remaining until the
deadline. The scheduling problem then is to determine an optimal
strategy.



\begin{proposition}
\label{prop:np}
The above scheduling problem is NP-hard, even for a one-PE, failure-free platform
with zero I/O costs, two modes per task, and
zero restart overhead.
\end{proposition}

\begin{proof}
With no failures, the execution is non-stochastic. All scheduling decisions
can be made at time zero and, with one PE, we can restrict ourselves to
strategies that execute each task exactly once with one mode.  Let us
prove the NP-completeness of the decision problem: does there exist a
strategy that can achieve an output error of at most some threshold $\tau$?
The problem is in NP because, given the list of the chosen execution modes,
one can compute the overall execution time and output error in polynomial
time.  We prove NP-completeness via a reduction from 0-1
\knapsack~\cite{GareyJ79}.

Consider an instance of 0-1 \knapsack with $m$ items, where item
$i{=}1,\ldots,m$ has weight $w_i$ and value $v_i$. The knapsack capacity is
$W$, and for $V{>}0$ we consider the problem of deciding whether a solution
exists with total value ${\geq}V$.  We construct an instance
of our scheduling decision problem as follows. We consider an $m$-task
chain in which task $T_i$, $i{=}1,\ldots,m$ has two execution modes:
(i)~\emph{fast}: execution time $t(x,y){=}1$, output size $d(x,y){=}0$, 
and error $e(x,y){=}y{+}v_i$ (i.e., introduced additive error); and (ii)~\emph{slow}: execution
	time $t(x,y){=}1{+}w_i$, output size $d(x,y){=}0$, and error
function $e(x,y){=}y$ (i.e., no introduced error).
We set the deadline to $D{=}m{+}W$, the size of the initial input to 0,
and the error of the initial input to 1. 
Let us define $\tau{=}1{+}V_{\Sigma}{-}V$, where
$V_{\Sigma}{=}\sum_{i=1}^m v_i$.  We consider the problem of deciding whether
there exists a strategy for which the chain's output error is
${\leq}\tau$.  We define $e_f{=}2{+}V_{\Sigma}$, so that a chain execution
that does not complete by the deadline leads to error strictly greater than
$\tau$.  This construction is computable in polynomial time.

Let $S$ be the set of tasks that are executed with their slow execution
mode, as decided by the strategy. 
The application execution finishes by deadline $D$ when:
$$
\sum_{i \notin S} 1 + \sum_{i \in S} (1 + w_i) \leq m + W\;,
\text{i.e., when}
\;\sum_{i \in S} w_i \leq W .
$$
The chain's output error is:
$$
1 + \sum_{i \notin S} v_i = 1 + V_{\Sigma} - \sum_{i \in S} v_i = \tau + V - \sum_{i \in S} v_i \;.
$$
Therefore, the chain's output error is ${\leq}\tau$ if and only if $\sum_{i
\in S} v_i{\geq}V$.  In other words, the answer to the scheduling decision
problem is yes if and only if the answer to the knapsack decision problem is yes,
completing the reduction.
\end{proof}



